% appendix_new.tex : theory appendix for LRC-BART (locally robust commensurate BART).
% Standalone body: compile by \input-ing this file into a wrapper that carries the
% preamble of paper/main.tex (see 03-theory/build/wrapper.tex). No \documentclass here.
%
% Bibliography keys used below exist in paper/ref.bib except the two flagged in
% comments (Hahn, Murray and Carvalho 2020; Hobbs et al. 2011), which are cited in
% text only until the entries are added.

% ---- Appendix-local theorem environments (Theorem A1, Lemma A1, ...) -------------
\newtheorem{thmA}{Theorem}
\renewcommand{\thethmA}{A\arabic{thmA}}
\newtheorem{lemA}{Lemma}
\renewcommand{\thelemA}{A\arabic{lemA}}
\newtheorem{propA}{Proposition}
\renewcommand{\thepropA}{A\arabic{propA}}
\newtheorem{corA}{Corollary}
\renewcommand{\thecorA}{A\arabic{corA}}
\theoremstyle{definition}
\newtheorem{assA}{Assumption}
\renewcommand{\theassA}{A\arabic{assA}}
\newtheorem{remA}{Remark}
\renewcommand{\theremA}{A\arabic{remA}}
\theoremstyle{plain}
\renewcommand{\theequation}{A\arabic{equation}}
\newcommand{\Pspike}{P_{\mathrm{spike}}}
\newcommand{\ess}{\mathrm{ESS}}
\newcommand{\Phibar}{\bar\Phi}
\newcommand{\tr}{^{\top}}

\section*{Web Appendix A: Model, notation, and exact leaf marginals}

This appendix collects the proofs for the locally robust commensurate BART (LRC-BART) model of the main text. We fix notation first. The two control sources are indexed by $g=1$ (RCT control) and $g=2$ (RWD control), with residual variances $\sigma_1^2$ and $\sigma_2^2$. The control outcomes follow
\begin{equation}
y_i=f(x_i)+D_i\,g(x_i)+\varepsilon_i,\qquad \varepsilon_i\sim\Normal(0,\sigma^2_{g_i}),\qquad D_i=\mathbf 1\{i\ \text{is an RCT control}\},
\label{eqA:model}
\end{equation}
where $f$ is a standard BART ensemble of $H_f$ trees with leaf prior $\Normal(0,\lambda_f^2)$ and $g$ is an ensemble of $H_g$ trees whose leaves carry the two-component prior
\begin{equation}
\theta_{h\ell}\given z_{h\ell},\tau_0^2,\tau_1^2\sim\Normal\!\left(0,\tau_{1-z_{h\ell}}^2\right),\qquad z_{h\ell}\given w\iid\Ber(w),
\label{eqA:leafprior}
\end{equation}
with $z=1$ the spike (commensurate component) and $z=0$ the slab, $w\sim\mathrm{Beta}(a_w,b_w)$, $\tau_1^2$ fixed by the range rule, and
\begin{equation}
\tau_0^2\sim\text{scaled-Inv-}\chi^2(\nu_0,s_0^2)\ \text{truncated to }(0,\tau_1^2),
\label{eq:tau0prior}
\end{equation}
which is the prior the sampler implements (its conditional for $\tau_0^2$ is the scaled-Inv-$\chi^2(\nu_0+K_0,\cdot)$ law truncated to $(0,\tau_1^2)$, which is conjugate for \eqref{eq:tau0prior} and not for the untruncated prior). Throughout we impose
\begin{assA}[(A1)]\label{ass:A1}
$0<\tau_0^2<\tau_1^2<\infty$.
\end{assA}
\noindent Under \eqref{eq:tau0prior}, Assumption~\ref{ass:A1} holds almost surely under the prior and along the chain. The variance-based calibration in Web Appendix~D uses the untruncated scaled-Inv-$\chi^2(\nu_0,s_0^2)$ law. The mass removed by the fitting prior is $P(\chi^2_{\nu_0}<\nu_0s_0^2/\tau_1^2)$ and depends on the fitted calibration scale; the two laws are distinguished in that appendix.

Within a leaf, $n_g$ and $\bar y_g$ denote the number and the mean of the source-$g$ partial residuals routed to the leaf, $v_g=\sigma_g^2/n_g$, and $\phi(\cdot;m,V)$ is the $\Normal(m,V)$ density; $\Phi$ is the standard normal distribution function and $\Phibar=1-\Phi$. For a $g$ leaf only RCT controls enter, and we write $\bar y=\bar y_1$, $v=v_1$, $n=n_1$ when no confusion arises. The identity we use repeatedly is the Gaussian product rule
\begin{equation}
\phi(\theta;0,\tau^2)\,\phi(\bar y;\theta,v)=\phi(\bar y;0,\tau^2+v)\,\phi\Bigl(\theta;\tfrac{\tau^2}{\tau^2+v}\bar y,\ \tfrac{\tau^2 v}{\tau^2+v}\Bigr),
\label{eq:gaussprod}
\end{equation}
which follows by completing the square in $\theta$.

\subsection*{A.1 Exact leaf marginals}

\begin{lemA}[Exact leaf marginals]\label{lem:marg}
Fix all trees other than tree $h$, their leaf parameters, the residual variances, $\lambda_f^2$, and $(\tau_0^2,\tau_1^2,w)$.
\begin{enumerate}[label=(\alph*),leftmargin=2.2em]
\item (\emph{$f$ leaf.}) Let $r_i=y_i-f_{-h}(x_i)-D_ig(x_i)$ be the partial residuals of the control observations routed to a leaf of $f$-tree $h$, so that $r_i\given\theta\indep\Normal(\theta,\sigma^2_{g_i})$ and $\theta\sim\Normal(0,\lambda_f^2)$. With $A=n_1/\sigma_1^2+n_2/\sigma_2^2$ and $B=S_1/\sigma_1^2+S_2/\sigma_2^2$, $S_g=\sum_{i:g_i=g}r_i$,
\begin{equation}
\int\prod_i\phi(r_i;\theta,\sigma^2_{g_i})\,\phi(\theta;0,\lambda_f^2)\,\dd\theta
=\Bigl[\prod_i\phi(r_i;0,\sigma^2_{g_i})\Bigr]\,M_f,\qquad
M_f=\frac{\exp\bigl\{\lambda_f^2B^2/\{2(1+\lambda_f^2A)\}\bigr\}}{(1+\lambda_f^2A)^{1/2}}.
\label{eqA:Mf}
\end{equation}
\item (\emph{$g$ leaf.}) Let $r_i=y_i-f(x_i)-g_{-h}(x_i)$ be the partial residuals of the $n_1$ RCT controls routed to a leaf of $g$-tree $h$, with mean $\bar y$. Under the prior \eqref{eqA:leafprior},
\begin{equation}
\sum_{z\in\{0,1\}}P(z)\int\prod_i\phi(r_i;\theta,\sigma_1^2)\,\phi(\theta;0,\tau_{1-z}^2)\,\dd\theta
=\Bigl[\prod_i\phi(r_i;0,\sigma_1^2)\Bigr]\,M_g,
\label{eqA:Mg}
\end{equation}
\begin{equation}
M_g=\frac{w\,\phi(\bar y;0,\tau_0^2+v)+(1-w)\,\phi(\bar y;0,\tau_1^2+v)}{\phi(\bar y;0,v)},
\label{eq:Mg2}
\end{equation}
where the index $1-z$ selects $\tau_0^2$ for $z=1$ and $\tau_1^2$ for $z=0$. A leaf with $n_1=0$ has $M_g=1$.
\item (\emph{Leaf conditionals.}) In a $g$ leaf,
\begin{equation}
\begin{aligned}
P(z=1\given\bar y)&=\frac{w\,\phi(\bar y;0,\tau_0^2+v)}{w\,\phi(\bar y;0,\tau_0^2+v)+(1-w)\,\phi(\bar y;0,\tau_1^2+v)},\\
\theta\given z,\bar y&\sim\Normal\Bigl(\frac{\tau_{1-z}^2}{\tau_{1-z}^2+v}\bar y,\ \frac{\tau_{1-z}^2v}{\tau_{1-z}^2+v}\Bigr).
\end{aligned}
\label{eq:leafcond}
\end{equation}
and in an $f$ leaf $\theta\given\cdot\sim\Normal\bigl(\lambda_f^2B/(1+\lambda_f^2A),\ \lambda_f^2/(1+\lambda_f^2A)\bigr)$.
\end{enumerate}
\end{lemA}

\begin{proof}
(a) Expanding $-(r_i-\theta)^2/(2\sigma^2_{g_i})=-r_i^2/(2\sigma^2_{g_i})+r_i\theta/\sigma^2_{g_i}-\theta^2/(2\sigma^2_{g_i})$ and summing over $i$ gives
$\prod_i\phi(r_i;\theta,\sigma^2_{g_i})=\prod_i\phi(r_i;0,\sigma^2_{g_i})\exp(\theta B-\theta^2A/2)$. Multiplying by $\phi(\theta;0,\lambda_f^2)$ and integrating, the exponent is $-\tfrac12(A+\lambda_f^{-2})\theta^2+B\theta$, a Gaussian integral equal to $(2\pi)^{1/2}(A+\lambda_f^{-2})^{-1/2}\exp\{B^2/(2(A+\lambda_f^{-2}))\}$; dividing by the prior normalizing constant $(2\pi\lambda_f^2)^{1/2}$ gives $M_f$. The conditional of $\theta$ is read off the same quadratic form.

(b) With a single source, (a) gives $\prod_i\phi(r_i;\theta,\sigma_1^2)=\prod_i\phi(r_i;0,\sigma_1^2)\exp\{\theta n\bar y/\sigma_1^2-\theta^2n/(2\sigma_1^2)\}$, and the exponential equals $\phi(\bar y;\theta,v)/\phi(\bar y;0,v)$, since $-\{(\bar y-\theta)^2-\bar y^2\}/(2v)=(2\bar y\theta-\theta^2)n/(2\sigma_1^2)$. Hence the likelihood factors as $\prod_i\phi(r_i;0,\sigma_1^2)\times\phi(\bar y;\theta,v)/\phi(\bar y;0,v)$, the sufficiency of $\bar y$. Integrating $\phi(\bar y;\theta,v)$ against $\phi(\theta;0,\tau_{1-z}^2)$ by \eqref{eq:gaussprod} gives $\phi(\bar y;0,\tau_{1-z}^2+v)$, and the sum over $z$ with weights $w$ and $1-w$ gives $M_g$. When $n_1=0$ the leaf has no likelihood factor and the integral of the prior is one.

(c) By \eqref{eq:gaussprod}, the joint density of $(z,\theta)$ given $\bar y$ is proportional to
\[
P(z)\phi(\bar y;0,\tau_{1-z}^2+v)\,
\phi\Bigl(\theta;\frac{\tau_{1-z}^2}{\tau_{1-z}^2+v}\bar y,\frac{\tau_{1-z}^2v}{\tau_{1-z}^2+v}\Bigr).
\]
Integrating out $\theta$ gives the stated Bernoulli weight; the remaining factor is the conditional density of $\theta$.
\end{proof}

\begin{propA}[The tree moves target the exact posterior]\label{prop:exact}
Condition on $(\sigma_1^2,\sigma_2^2,\lambda_f^2,\tau_0^2,\tau_1^2,w)$ and on all trees other than $h$ together with their leaf parameters. Let the update of tree $h$ consist of a Metropolis--Hastings step on the partition $T_h$ with grow, prune, and change proposals, acceptance ratio equal to the tree-prior ratio times the proposal ratio times the ratio of products of $M_f$ (for $f$ trees) or $M_g$ (for $g$ trees) over the affected leaves, followed by a draw of the leaf parameters of $T_h$ (and, for a $g$ tree, the indicators) from \eqref{eq:leafcond}. Then the composite update leaves the conditional posterior of $(T_h,\theta_h)$ (resp.\ $(T_h,\theta_h,z_h)$) invariant. Consequently the full sweep, with the conjugate updates of $(\sigma_g^2,\tau_0^2,w)$, the update of $\tau_0^2$ being the truncated conditional under the prior \eqref{eq:tau0prior}, is a partially collapsed Gibbs sampler whose invariant law is the exact joint posterior under \eqref{eqA:leafprior} and \eqref{eq:tau0prior}.
\end{propA}

\begin{proof}
Given the conditioning, the marginal conditional of $T_h$ with the leaf parameters and indicators integrated out is $\pi(T_h)\prod_i\phi(r_i;0,\sigma^2_{g_i})\prod_{\ell\in T_h}M_\ell$ by Lemma~\ref{lem:marg}, where $M_\ell$ is $M_f$ or $M_g$. The partial residuals $r_i$ do not depend on $T_h$, and the observation-level product is over all observations of the tree, so it is the same for the current and the proposed partition and cancels from the Metropolis--Hastings ratio. The stated acceptance ratio is therefore the exact ratio for the collapsed target; for a grow move that splits leaf $P$ into $L$ and $R$, the leaf factor ratio is $M_LM_R/M_P$, and the prune and change ratios follow in the same way. A Metropolis--Hastings kernel with the exact ratio leaves the collapsed conditional invariant. Following that step by a draw of the leaf parameters (and indicators) from their exact conditional given $T_h$ gives a kernel that preserves the joint conditional of $(T_h,\theta_h,z_h)$, which is the marginal-then-conditional construction of \citet{Chipman2010BART}. The conjugate updates of $\sigma_g^2$, $\tau_0^2$, and $w$ are ordinary Gibbs steps, so the sweep is a partially collapsed Gibbs sampler with the joint posterior as invariant law.
\end{proof}

\begin{remA}[No surrogate is needed]\label{rem:nosurrogate}
Every factor entering a tree move is closed form because, conditionally on $(\tau_0^2,\tau_1^2,w)$, the leaf prior \eqref{eqA:leafprior} is a finite Gaussian mixture, and the Gaussian likelihood integrates against each component in closed form. If a leaf-specific discrepancy variance instead has a scaled-Inv-$\chi^2$ prior and is integrated inside the leaf marginal, the resulting additive scale mixture is $\int\phi(\bar y_1;0,v_1+\tau^2)\phi(\bar y_2;0,v_2+\tau^2)p(\tau^2)\,\dd\tau^2$, with no elementary form. In LRC-BART the scale $\tau_0^2$ is a global parameter updated by its own conjugate step, and the local adaptivity is carried by the per-leaf indicator $z$, so the leaf marginal is exact and the tree kernels target the exact posterior (Proposition~\ref{prop:exact}). Posterior invariance therefore follows from the exact update construction without a surrogate-likelihood perturbation bound; Web Appendix~F states its scope.
\end{remA}

\section*{Web Appendix B: Bounded local influence (T1)}

\subsection*{B.1 A single $g$ leaf}

We fix a $g$ leaf and condition on $f$, on the other $g$ trees, on the partition, and on $(\sigma_1^2,\tau_0^2,\tau_1^2,w)$, so that $\bar y\given\theta\sim\Normal(\theta,v)$ with $v=\sigma_1^2/n$. Write
\begin{equation}
\rho_j=\frac{\tau_j^2}{\tau_j^2+v},\qquad \hat\theta_j=\rho_j\,\bar y,\qquad V_j=\rho_j\,v=\frac{\tau_j^2v}{\tau_j^2+v},\qquad j\in\{0,1\},
\label{eq:rho}
\end{equation}
so that $\hat\theta_1$ is the slab estimate (close to the RCT-only residual mean $\bar y$ when $\tau_1^2\gg v$) and $\hat\theta_0$ the spike-shrunk estimate (close to zero, that is, to the pooled surface $f$, when $\tau_0^2\ll v$). The subscript $j$ of $\rho_j$, $\hat\theta_j$, $V_j$ indexes the scale $\tau_j^2$, so $j=1$ is the slab, whereas the indicator convention of \eqref{eqA:leafprior} has $z=1$ for the spike. Under Assumption~\ref{ass:A1}, $0<\rho_0<\rho_1<1$. Let $p=P(z=0\given\bar y)$ be the posterior slab probability and define
\begin{equation}
\kappa=\frac12\Bigl(\frac{1}{\tau_0^2+v}-\frac{1}{\tau_1^2+v}\Bigr),\qquad
a=\log\frac{1-w}{w}+\frac12\log\frac{\tau_0^2+v}{\tau_1^2+v},\qquad
C=e^{-a}=\frac{w}{1-w}\sqrt{\frac{\tau_1^2+v}{\tau_0^2+v}}.
\label{eq:kappaC}
\end{equation}
Under Assumption~\ref{ass:A1}, $\kappa>0$, and a direct computation gives
\begin{equation}
\rho_1-\rho_0=\frac{v(\tau_1^2-\tau_0^2)}{(\tau_1^2+v)(\tau_0^2+v)}=2v\kappa .
\label{eq:rhodiff}
\end{equation}

\begin{thmA}[Bounded local influence]\label{thm:T1}
Under Assumption~\ref{ass:A1} and the conditioning above, the following hold for every $\bar y\in\mathbb R$; parts (b) to (d) assume $0<w<1$.
\begin{enumerate}[label=(\alph*),leftmargin=2.2em]
\item (\emph{Posterior.}) $\theta\given\bar y\sim p\,\Normal(\hat\theta_1,V_1)+(1-p)\,\Normal(\hat\theta_0,V_0)$, and hence
\begin{equation}
\E[\theta\given\bar y]=p\,\frac{\tau_1^2}{\tau_1^2+v}\,\bar y+(1-p)\,\frac{\tau_0^2}{\tau_0^2+v}\,\bar y .
\label{eq:postmean}
\end{equation}
\item (\emph{Slab probability.}) $\displaystyle \logit p=\log\frac{p}{1-p}=a+\kappa\,\bar y^2$, with $a$ and $\kappa$ as in \eqref{eq:kappaC}.
\item (\emph{Bounded spike contribution and bounded RWD-induced shift.}) The spike-shrunk part of \eqref{eq:postmean} satisfies $|(1-p)\hat\theta_0|\le|\bar y|\,\tau_0^2/(\tau_0^2+v)$. The shift of the posterior mean relative to the slab estimate,
\begin{equation}
I(\bar y)=\E[\theta\given\bar y]-\hat\theta_1=-(1-p)\,(\rho_1-\rho_0)\,\bar y ,
\label{eq:I}
\end{equation}
satisfies $|I(\bar y)|\le(\rho_1-\rho_0)\,|\bar y|\,\min\{1,\,Ce^{-\kappa\bar y^2}\}$ and
\begin{equation}
\sup_{\bar y\in\mathbb R}|I(\bar y)|\ \le\ 2v\sqrt{\kappa\max\{\log C,\tfrac12\}}\ \le\ \sqrt{2v\max\{\log C,\tfrac12\}} .
\label{eq:supI}
\end{equation}
\item (\emph{Gaussian rate.}) $1-p\le Ce^{-\kappa\bar y^2}$, so $p\to1$ and $I(\bar y)\to0$ as $|\bar y|\to\infty$, at a Gaussian rate.
\item (\emph{Nested cases.}) With $w=1$ (C-BART), $p\equiv0$ and $I(\bar y)=-(\rho_1-\rho_0)\bar y$, linear and unbounded; if in addition $\tau_0^2\to0$ (BART-CP), $\E[\theta\given\bar y]\to0$ and $I(\bar y)\to-\rho_1\bar y$. With $w=0$, $p\equiv1$ and $I\equiv0$. If Assumption~\ref{ass:A1} fails with $\tau_0^2>\tau_1^2$, then $\kappa<0$, the roles of the two components are exchanged, and (c), (d) hold with the wide component in the role of the slab; if $\tau_0^2=\tau_1^2$, $p\equiv1-w$ and $I\equiv0$.
\end{enumerate}
\end{thmA}

\begin{proof}
(a) The posterior density is proportional to $\{w\phi(\theta;0,\tau_0^2)+(1-w)\phi(\theta;0,\tau_1^2)\}\phi(\bar y;\theta,v)$. Applying \eqref{eq:gaussprod} to each term gives $w\phi(\bar y;0,\tau_0^2+v)\phi(\theta;\hat\theta_0,V_0)+(1-w)\phi(\bar y;0,\tau_1^2+v)\phi(\theta;\hat\theta_1,V_1)$, a two-component Gaussian mixture with weights proportional to $w\phi(\bar y;0,\tau_0^2+v)$ and $(1-w)\phi(\bar y;0,\tau_1^2+v)$; normalizing gives $1-p$ and $p$ as in \eqref{eq:leafcond}, and \eqref{eq:postmean} is the mean of the mixture.

(b) The log odds of the two weights are
\[
\log\frac{(1-w)\phi(\bar y;0,\tau_1^2+v)}{w\,\phi(\bar y;0,\tau_0^2+v)}
=\log\frac{1-w}{w}-\frac12\log\frac{\tau_1^2+v}{\tau_0^2+v}+\frac{\bar y^2}{2}\Bigl(\frac{1}{\tau_0^2+v}-\frac{1}{\tau_1^2+v}\Bigr)=a+\kappa\bar y^2 .
\]

(c) The first bound is immediate from $0\le1-p\le1$. For $I$, subtract $\hat\theta_1$ from \eqref{eq:postmean}. Since $1-p=(1+e^{a+\kappa\bar y^2})^{-1}\le\min\{1,e^{-a-\kappa\bar y^2}\}=\min\{1,Ce^{-\kappa\bar y^2}\}$, the pointwise bound follows. For the supremum we show that for $\kappa>0$ and $C>0$,
\begin{equation}
\sup_{u\in\mathbb R}|u|\min\{1,Ce^{-\kappa u^2}\}\le\sqrt{\max\{\log C,\tfrac12\}/\kappa}.
\label{eq:sup-lemma}
\end{equation}
Put $t=\kappa u^2$. On $\{t\le\log C\}$ (empty if $C\le1$) the left side is at most $|u|\le\sqrt{\log C/\kappa}$. On $\{t>\max(\log C,0)\}$ the left side equals $\sqrt{t/\kappa}\,Ce^{-t}$. If $\log C\ge\tfrac12$, the function $\sqrt t\,e^{-t}$ is decreasing on $[\tfrac12,\infty)$, so the value is at most $\sqrt{\log C/\kappa}\,Ce^{-\log C}=\sqrt{\log C/\kappa}$. If $\log C<\tfrac12$, then $\sup_{t>0}\sqrt t\,e^{-t}=(2e)^{-1/2}$ gives at most $C(2e\kappa)^{-1/2}<\sqrt{1/(2\kappa)}$ because $C<e^{1/2}$. This proves \eqref{eq:sup-lemma}. Combining with \eqref{eq:rhodiff}, $\sup|I|\le2v\kappa\sqrt{\max\{\log C,\tfrac12\}/\kappa}=2v\sqrt{\kappa\max\{\log C,\tfrac12\}}$, and $\kappa\le1/(2(\tau_0^2+v))\le1/(2v)$ gives the second inequality in \eqref{eq:supI}.

(d) is the pointwise bound in (c) together with $\kappa>0$. (e) With $w=1$ the slab weight is zero for every $\bar y$; with $w=0$ the spike weight is zero; the remaining statements are read off \eqref{eq:kappaC} and \eqref{eq:I}.
\end{proof}

\begin{remA}[What is and is not bounded]\label{rem:slabshrink}
Theorem~\ref{thm:T1} bounds the influence of the commensurate component relative to the slab estimate $\hat\theta_1=\rho_1\bar y$. The slab estimate itself is shrunk toward zero, that is toward the pooled surface $f$, by the factor $\rho_1$: $\E[\theta\given\bar y]-\bar y=-(1-\rho_1)\bar y+I(\bar y)$ with $1-\rho_1=v/(\tau_1^2+v)$. This linear term is the range-rule regularization that any BART leaf with prior variance $\tau_1^2$ applies, and it is of the same form as the shrinkage of a BART-NP leaf toward its ensemble center; at $n=50$, $\sigma_1=1.5$ and $\tau_1^2=1.25$ its slope is $0.035$. The bounded term $I$ is the part specific to borrowing.
\end{remA}

\begin{remA}[Numerical illustration]\label{rem:T1num}
With $\sigma_1=1.5$, $n=50$, $\tau_0^2=0.001$ (the scale calibrated for $N_{\rm target}=100$ in the Gaussian study), $\tau_1^2=1.25$ (the range rule at $H_g=5$ for an outcome range of $10$) and $w=0.8$ fixed for illustration, $v=0.045$, $\kappa=10.5$, $C=21.2$, $\log C=3.06$, and $\rho_1-\rho_0=0.94$. The bound \eqref{eq:supI} gives $\sup|I|\le0.51$; the exact supremum of \eqref{eq:I}, a one-dimensional maximization, is $0.31$, attained at $\bar y=0.43$. By comparison C-BART's shift is $0.94\,|\bar y|$, which exceeds $0.51$ once $|\bar y|>0.54$.
\end{remA}

\subsection*{B.2 The estimand level: fixed partitions of the $g$ ensemble}

The estimand is the RCT-standardized control mean $\mu=N^{-1}\sum_{i\le N}f_1(x_i)$ over the $N$ RCT covariate profiles, and $f_1=f+g$, so $\mu=\mu_f+\mu_g$ with $\mu_g=N^{-1}\sum_{i\le N}g(x_i)$. We condition on $f$, on the partitions $\mathcal T=(T_1,\dots,T_{H_g})$ of the $g$ trees, and on $(\sigma_1^2,\tau_0^2,\tau_1^2,w)$; the RWD then enters $\mu_g$ only through the shrinkage of $g$ toward zero. Index the $L_g$ leaves of all $g$ trees by $j$, let $\theta\in\mathbb R^{L_g}$ collect the leaf parameters and $z\in\{0,1\}^{L_g}$ the indicators, and let $r\in\mathbb R^{n_1}$ be the residuals $y_i-f(x_i)$ of the RCT controls. Let $X\in\{0,1\}^{n_1\times L_g}$ be the incidence matrix ($X_{ij}=1$ if control $i$ falls in leaf $j$), so each row of $X$ has exactly $H_g$ ones and $g(x_i)=(X\theta)_i$. Let $a\in\mathbb R^{L_g}$ with $a_j=N_j/N$, $N_j$ the number of estimand profiles in leaf $j$, so $\mu_g=a\tr\theta$, $\|a\|_1=H_g$, and $\|a\|_2^2\le\max_ja_j\sum_ja_j\le H_g$, so $\|a\|_2\le\sqrt{H_g}$. Then
\begin{equation}
r\given\theta\sim\Normal(X\theta,\sigma_1^2I),\qquad \theta\given z\sim\Normal(0,D_z),\quad D_z=\mathrm{diag}(\tau^2_{1-z_j}),\qquad \pi(z)=\prod_jw^{z_j}(1-w)^{1-z_j}.
\label{eq:linform}
\end{equation}
Standard Gaussian algebra gives, with $V_z=XD_zX\tr+\sigma_1^2I$ and $\Sigma_z=(X\tr X/\sigma_1^2+D_z^{-1})^{-1}$,
\begin{equation}
P(z\given r)\propto\pi(z)\,\phi_{n_1}(r;0,V_z),\qquad
\theta\given z,r\sim\Normal\bigl(\hat\theta(z),\Sigma_z\bigr),\qquad \hat\theta(z)=D_zX\tr V_z^{-1}r,
\label{eqA:zpost}
\end{equation}
where the form of $\hat\theta(z)$ follows from $(X\tr X/\sigma^2+D^{-1})^{-1}X\tr/\sigma^2=DX\tr(XDX\tr+\sigma^2I)^{-1}$. We write $\mathbf 0$ for the all-slab configuration; the $w=0$ model (the tree analogue of BART-PP, a free source offset regularized by the range rule) has $\E_{w=0}[\mu_g\given r]=a\tr\hat\theta(\mathbf 0)$, and we call
\begin{equation}
\Delta(r)=\E[\mu_g\given r]-a\tr\hat\theta(\mathbf 0)=\sum_{z\ne\mathbf 0}P(z\given r)\,a\tr\bigl(\hat\theta(z)-\hat\theta(\mathbf 0)\bigr)
\label{eq:Delta}
\end{equation}
the RWD-induced shift of the estimand at the given partitions.

\begin{corA}[Bounded shift of the standardized control mean]\label{cor:estimand}
Under Assumption~\ref{ass:A1}, $0<w<1$, and the conditioning of this subsection, let $\delta=\tau_1^2-\tau_0^2$ and $b=w\tau_1/\{(1-w)\tau_0\}$.
\begin{enumerate}[label=(\alph*),leftmargin=2.2em]
\item (\emph{General partitions.}) $\displaystyle\sup_{r\in\mathbb R^{n_1}}|\Delta(r)|\le \sqrt{H_g}\Bigl(1+\frac{\tau_1}{\tau_0}\Bigr)(2^{L_g}-1)\sqrt{2\delta\max\{L_g\log b,\tfrac12\}}<\infty$.
\item (\emph{Single tree, $H_g=1$.}) With leaves $\ell$, RCT counts $n_{1\ell}$, $v_\ell=\sigma_1^2/n_{1\ell}$, leaf means $\bar y_\ell$, and $\kappa_\ell,C_\ell$ as in \eqref{eq:kappaC} with $v=v_\ell$,
\[
\Delta(r)=\sum_{\ell:\,n_{1\ell}>0}a_\ell\,I_\ell(\bar y_\ell),\qquad
\sup_r|\Delta(r)|\le\sum_\ell a_\ell\sqrt{2v_\ell\max\{\log C_\ell,\tfrac12\}} ,
\]
and the right side is at most $\max_\ell\sqrt{2v_\ell\max\{\log C_\ell,\tfrac12\}}$.
\item (\emph{$H_g$ trees with a common partition.}) If all $g$ trees share one partition with leaves $\ell$, write $\vartheta_\ell=\sum_h\theta_{h\ell}$, $T_k=k\tau_0^2+(H_g-k)\tau_1^2$, $\rho^{(k)}_\ell=T_k/(T_k+v_\ell)$,
\[
\kappa_{k\ell}=\frac12\Bigl(\frac1{T_k+v_\ell}-\frac1{T_0+v_\ell}\Bigr),\qquad
C_{k\ell}=\binom{H_g}{k}\Bigl(\frac{w}{1-w}\Bigr)^{k}\sqrt{\frac{T_0+v_\ell}{T_k+v_\ell}} .
\]
Then $P(k\ \text{spikes in leaf }\ell\given\bar y_\ell)\propto\binom{H_g}{k}w^k(1-w)^{H_g-k}\phi(\bar y_\ell;0,T_k+v_\ell)$, the leaves are independent a posteriori,
\begin{align}
\Delta(r)&=-\sum_\ell a_\ell\,\bar y_\ell\sum_{k=1}^{H_g}P(k\given\bar y_\ell)\bigl(\rho^{(0)}_\ell-\rho^{(k)}_\ell\bigr),
\label{eq:Delta-common}\\
\sup_r|\Delta(r)|&\le\max_\ell\ 2v_\ell\sum_{k=1}^{H_g}\sqrt{\kappa_{k\ell}\max\{\log C_{k\ell},\tfrac12\}}
\ \le\ \max_\ell H_g\,v_\ell\sqrt{\frac{2\max_k\max\{\log C_{k\ell},\tfrac12\}}{H_g\tau_0^2+v_\ell}} .
\label{eq:Delta-common-bound}
\end{align}
\item (\emph{Contrasts.}) Under complete pooling of the control surface ($w=1$, $\tau_0^2\to0$) the shift is $-a\tr\hat\theta(\mathbf 0)$, and under C-BART ($w=1$) it is $a\tr(\hat\theta(\mathbf 1)-\hat\theta(\mathbf 0))$; both are linear in $r$, and each is unbounded on $\mathbb R^{n_1}$ unless the corresponding linear functional vanishes identically. For a single tree the C-BART functional is $-\sum_\ell a_\ell(\rho_{1\ell}-\rho_{0\ell})\bar y_\ell$, which is nonzero whenever some leaf has $a_\ell>0$ and $n_{1\ell}>0$.
\end{enumerate}
\end{corA}

\begin{proof}
(a) Fix $z\ne\mathbf 0$, let $S=\{j:z_j=1\}$ be its spike leaves, $E_S$ the diagonal indicator of $S$, $X_S$ the corresponding columns of $X$, and $u=X_S\tr V_{\mathbf 0}^{-1}r\in\mathbb R^{|S|}$. Since $D_{\mathbf 0}-D_z=\delta E_S$,
\[
V_z^{-1}-V_{\mathbf 0}^{-1}=V_z^{-1}(V_{\mathbf 0}-V_z)V_{\mathbf 0}^{-1}=\delta\,V_z^{-1}X_SX_S\tr V_{\mathbf 0}^{-1},
\qquad\text{so}\qquad
V_z^{-1}r=V_{\mathbf 0}^{-1}r+\delta\,V_z^{-1}X_Su .
\]
Three consequences follow. First, the quadratic form in the likelihood ratio is
\[
r\tr(V_z^{-1}-V_{\mathbf 0}^{-1})r=\delta\,(V_z^{-1}r)\tr X_Su=\delta\bigl(u+\delta X_S\tr V_z^{-1}X_Su\bigr)\tr u=\delta\|u\|^2+\delta^2u\tr X_S\tr V_z^{-1}X_Su\ \ge\ \delta\|u\|^2 .
\]
Second, the difference of posterior means is
\[
\hat\theta(z)-\hat\theta(\mathbf 0)=(D_z-D_{\mathbf 0})X\tr V_{\mathbf 0}^{-1}r+D_zX\tr(V_z^{-1}-V_{\mathbf 0}^{-1})r=-\delta\,\iota_Su+\delta\,D_zX\tr V_z^{-1}X_S\,u ,
\]
where $\iota_S$ embeds $\mathbb R^{|S|}$ into the $S$ coordinates of $\mathbb R^{L_g}$. The matrix $D_zX\tr V_z^{-1}X=D_z^{1/2}\{M\tr(MM\tr+\sigma_1^2I)^{-1}M\}D_z^{-1/2}$ with $M=XD_z^{1/2}$, and the bracket is symmetric with eigenvalues $\lambda_i/(\lambda_i+\sigma_1^2)\in[0,1)$, so $\|D_zX\tr V_z^{-1}X_S\|\le\|D_zX\tr V_z^{-1}X\|\le\|D_z^{1/2}\|\|D_z^{-1/2}\|\le\tau_1/\tau_0$, and $\|\hat\theta(z)-\hat\theta(\mathbf 0)\|\le\delta(1+\tau_1/\tau_0)\|u\|$. Third, $|V_z|=\sigma_1^{2n_1}|D_z|\,|D_z^{-1}+X\tr X/\sigma_1^2|$ and $D_z^{-1}\succeq D_{\mathbf 0}^{-1}$ give $|V_{\mathbf 0}|/|V_z|\le|D_{\mathbf 0}|/|D_z|=(\tau_1^2/\tau_0^2)^{|S|}$, so
\begin{align*}
P(z\given r)\le\frac{P(z\given r)}{P(\mathbf 0\given r)}
&=\frac{\pi(z)}{\pi(\mathbf 0)}\sqrt{\frac{|V_{\mathbf 0}|}{|V_z|}}\exp\Bigl\{-\tfrac12r\tr(V_z^{-1}-V_{\mathbf 0}^{-1})r\Bigr\}\\
&\le c_z\,e^{-\delta\|u\|^2/2},\qquad c_z=b^{|S|}.
\end{align*}
Combining, and using \eqref{eq:sup-lemma} in the variable $\|u\|$ with $\kappa=\delta/2$,
\begin{align*}
\bigl|P(z\given r)\,a\tr(\hat\theta(z)-\hat\theta(\mathbf 0))\bigr|
&\le\|a\|\,\delta\Bigl(1+\frac{\tau_1}{\tau_0}\Bigr)\|u\|\min\{1,c_ze^{-\delta\|u\|^2/2}\}\\
&\le \sqrt{H_g}\Bigl(1+\frac{\tau_1}{\tau_0}\Bigr)\sqrt{2\delta\max\{|S|\log b,\tfrac12\}} .
\end{align*}
Summing over the $2^{L_g}-1$ configurations $z\ne\mathbf 0$ and using $|S|\le L_g$ gives (a).

(b) For a single tree $X\tr X$ is diagonal, the residual vector splits by leaf, and \eqref{eq:linform} factorizes over leaves, so the leaves are independent a posteriori and each is the single-leaf model of Section~B.1 with $v=v_\ell$; leaves with $n_{1\ell}=0$ have posterior equal to prior under every $w$ and contribute zero to $\Delta$. Theorem~\ref{thm:T1}(c) bounds each $|I_\ell|$, and $\sum_\ell a_\ell=1$ gives the last inequality.

(c) With a common partition, the residual in leaf $\ell$ is $r_i=\vartheta_\ell+\varepsilon_i$ and, given the indicators, $\vartheta_\ell\sim\Normal(0,T_{k_\ell})$ with $k_\ell$ the number of spikes among the $H_g$ leaves stacked on $\ell$; the leaves are independent a posteriori and $\mu_g=\sum_\ell a_\ell\vartheta_\ell$. Summing $\pi(z)\phi(\bar y_\ell;0,T_k+v_\ell)$ over the $\binom{H_g}{k}$ configurations with $k$ spikes gives the stated posterior of $k$, and \eqref{eq:gaussprod} gives $\E[\vartheta_\ell\given k,\bar y_\ell]=\rho^{(k)}_\ell\bar y_\ell$. Under $w=0$, $k=0$ almost surely, which gives the expression for $\Delta$. For $k\ge1$, $P(k\given\bar y_\ell)\le P(k\given\bar y_\ell)/P(0\given\bar y_\ell)=C_{k\ell}e^{-\kappa_{k\ell}\bar y_\ell^2}$, and $\rho^{(0)}_\ell-\rho^{(k)}_\ell=v_\ell(T_0-T_k)/\{(T_0+v_\ell)(T_k+v_\ell)\}=2v_\ell\kappa_{k\ell}$; the $k$-th term is therefore at most $2v_\ell\kappa_{k\ell}|\bar y_\ell|\min\{1,C_{k\ell}e^{-\kappa_{k\ell}\bar y_\ell^2}\}\le2v_\ell\sqrt{\kappa_{k\ell}\max\{\log C_{k\ell},\tfrac12\}}$ by \eqref{eq:sup-lemma}. The final form uses $\kappa_{k\ell}\le\kappa_{H_g\ell}\le1/\{2(H_g\tau_0^2+v_\ell)\}$ and $\sum_\ell a_\ell=1$.

(d) With $w=1$ the posterior over $z$ is degenerate at $\mathbf 1$, so $\E[\mu_g\given r]=a\tr\hat\theta(\mathbf 1)$, linear in $r$; as $\tau_0^2\to0$, $\hat\theta(\mathbf 1)\to0$. A nonzero linear functional on $\mathbb R^{n_1}$ is unbounded. For a single tree, $\hat\theta_\ell(\mathbf 1)-\hat\theta_\ell(\mathbf 0)=(\rho_{0\ell}-\rho_{1\ell})\bar y_\ell$ by (b).
\end{proof}

\begin{remA}[On the constants]\label{rem:constants}
The constant in (a) is finite for every design and every partition, which is the content of the statement, but it is far from sharp: it treats each of the $2^{L_g}-1$ non-slab configurations separately and bounds the posterior probability of each by its odds against the all-slab configuration. Parts (b) and (c) are the cases in which the sum collapses, and they display the dependence on $H_g$, on $\tau_0$ through $\kappa$, and on the leaf shrinkage $\rho^{(H_g)}_\ell=H_g\tau_0^2/(H_g\tau_0^2+v_\ell)$ of the all-spike state. For the numbers of Remark~\ref{rem:T1num} with $\tau_0^2=0.0012$ and the study's configuration $H_g=5$, the common-partition bound in (c) is $1.15$ per leaf, while the exact supremum of \eqref{eq:Delta-common} is $0.148$, attained at $\bar y_\ell=0.33$ (at $H_g=10$ the corresponding values are $1.80$ and $0.067$ at $0.30$); the ensemble abstains sooner than a single leaf because there are $H_g$ ways to place one slab, and one slab leaf absorbs almost all of the discrepancy. The exact supremum is a one-dimensional maximization and is what we recommend reporting.
\end{remA}

\section*{Web Appendix C: Detection and the borrowing map (T2)}

\subsection*{C.1 A single $g$ leaf under a true local shift}

We keep the conditioning of Section~B.1 and add a data-generating assumption: the RCT residuals in the leaf are $r_i=d+\epsilon_i$ with $\epsilon_i\iid\Normal(0,\sigma_1^2)$, where $d$ is the true local discrepancy between the RCT control surface and the pooled surface $f$ (taken as fixed; see the scope discussion in Section~C.3 for the estimated $f$). Then $\bar y\sim\Normal(d,v)$, $v=\sigma_1^2/n$. Let $\Pspike(\bar y)=1-p=P(z=1\given\bar y)$.

\begin{thmA}[Leaf-level detection]\label{thm:T2}
Under Assumption~\ref{ass:A1} and the conditioning above:
\begin{enumerate}[label=(\alph*),leftmargin=2.2em]
\item (\emph{Spike probability.}) $\Pspike(\bar y)=\{1+\exp(a+\kappa\bar y^2)\}^{-1}$ with $a,\kappa$ as in \eqref{eq:kappaC}; it depends on $(\bar y,n)$ only through $\bar y^2$ and $v=\sigma_1^2/n$, is decreasing in $|\bar y|$, and satisfies $\Pspike(0)=\{1+e^{a}\}^{-1}$.
\item (\emph{Half-detection discrepancy.}) If $\log\frac{w}{1-w}+\frac12\log\frac{\tau_1^2+v}{\tau_0^2+v}>0$, the equation $\Pspike(\bar y)=\tfrac12$ has the unique positive solution
\begin{equation}
d_{1/2}=\sqrt{\frac{2(\tau_0^2+v)(\tau_1^2+v)}{\tau_1^2-\tau_0^2}\Bigl\{\log\frac{w}{1-w}+\frac12\log\frac{\tau_1^2+v}{\tau_0^2+v}\Bigr\}}\ ;
\label{eq:dhalf}
\end{equation}
otherwise $\Pspike\le\tfrac12$ for every $\bar y$ (with equality only at $\bar y=0$ when the bracket vanishes) and we set $d_{1/2}=0$. When $\tau_1^2\gg\tau_0^2+v$, $d_{1/2}^2\approx2(\tau_0^2+v)\{\log\frac{w}{1-w}+\frac12\log\frac{\tau_1^2}{\tau_0^2+v}\}$, a few multiples of $\tau_0^2+\sigma_1^2/n$. The probability that the leaf's slab probability exceeds one half, under the true shift $d$, is
\begin{equation}
P_d\bigl(p(\bar y)>\tfrac12\bigr)=P_d(|\bar y|>d_{1/2})=\Phibar\Bigl(\frac{d_{1/2}-d}{\sqrt v}\Bigr)+\Phibar\Bigl(\frac{d_{1/2}+d}{\sqrt v}\Bigr).
\label{eq:detectprob}
\end{equation}
\item (\emph{Large $n$, $\tau_0^2$ fixed.}) As $n\to\infty$, almost surely
\begin{equation}
\Pspike(\bar y)\ \to\ \Pspike^\infty(d)=\Bigl\{1+\frac{(1-w)\tau_0}{w\,\tau_1}\exp\Bigl(\frac{d^2}{2}\Bigl(\frac1{\tau_0^2}-\frac1{\tau_1^2}\Bigr)\Bigr)\Bigr\}^{-1}\in(0,1).
\label{eq:pinf}
\end{equation}
In particular, when $d=0$ (borrowing is warranted), $\Pspike\to w\tau_1/\{w\tau_1+(1-w)\tau_0\}>w$; when $d\ne0$, $\Pspike^\infty(d)\le\frac{w\tau_1}{(1-w)\tau_0}\exp\{-\frac{d^2}{2}(\tau_0^{-2}-\tau_1^{-2})\}$, which tends to zero as $|d|/\tau_0\to\infty$ but is strictly positive for every fixed $d$.
\item (\emph{Large $n$ with a vanishing spike.}) If $\tau_0^2=\tau_{0,n}^2\to0$ as $n\to\infty$ (with $\tau_1^2$, $w$ fixed), then for every $d\ne0$, $\Pspike(\bar y)\to0$ almost surely.
\end{enumerate}
\end{thmA}

\begin{proof}
(a) is Theorem~\ref{thm:T1}(b) with $\kappa>0$. (b) Setting $a+\kappa\bar y^2=0$ gives $\bar y^2=-a/\kappa$, which is \eqref{eq:dhalf} after substituting \eqref{eq:kappaC}; the condition is $-a>0$. The approximation drops $v$ and $\tau_0^2$ against $\tau_1^2$ in $(\tau_1^2+v)/(\tau_1^2-\tau_0^2)$. Since $p>\tfrac12$ iff $|\bar y|>d_{1/2}$ and $\bar y\sim\Normal(d,v)$, \eqref{eq:detectprob} follows.

(c) By the strong law, $\bar y\to d$ almost surely; $v\to0$ gives $a\to a_\infty=\log\frac{1-w}{w}+\frac12\log\frac{\tau_0^2}{\tau_1^2}$ and $\kappa\to\kappa_\infty=\frac12(\tau_0^{-2}-\tau_1^{-2})$, both finite. Hence $a+\kappa\bar y^2\to a_\infty+\kappa_\infty d^2$ almost surely, and \eqref{eq:pinf} follows by continuity of $t\mapsto(1+e^t)^{-1}$; the limit is strictly between $0$ and $1$ because the exponent is finite. The bound on $\Pspike^\infty(d)$ is $(1+e^t)^{-1}\le e^{-t}$.

(d) Write $\epsilon_n=\tau_{0,n}^2+v_n\to0$. Then $a_n+\kappa_n\bar y_n^2\ge\log\frac{1-w}{w}+\frac12\log\frac{\epsilon_n}{\tau_1^2+v_n}+\bar y_n^2\bigl(\frac1{2\epsilon_n}-\frac1{2\tau_1^2}\bigr)$. Almost surely $\bar y_n^2\to d^2>0$, and $\frac{d^2}{2\epsilon_n}+\frac12\log\epsilon_n\to\infty$, so the log odds of the slab diverge and $\Pspike\to0$.
\end{proof}

\paragraph{Fixed and vanishing spike scales.}
Theorem~\ref{thm:T2}(c) distinguishes increasing local sample size from decreasing spike scale: because the spike is a $\Normal(0,\tau_0^2)$ component and not a point mass, a discrepancy of the order of $\tau_0$ is consistent with it, and the limiting spike probability \eqref{eq:pinf} is positive for every $d$. What is true is (i) the limit is exponentially small in $d^2/(2\tau_0^2)$, so abstention is consistent in the regime $|d|/\tau_0\to\infty$, and (ii) abstention is consistent for every fixed $d\ne0$ when the spike scale shrinks with $n$ (part (d)). For a fixed-scale illustration, at $\tau_0^2=0.001$ a shift of $d=0.5$ has $d^2/(2\tau_0^2)=125$. The fitted two-arm sampler learns $\tau_0^2$ from the spike leaves under a prior calibrated to the order $10^{-3}$; the fixed-scale limit therefore does not establish its finite-sample detection behavior, which is assessed in Section~\ref{sec:sim}. The consistency statement that does hold for every fixed $d\ne0$, without a vanishing spike, is the one for the discrepancy map in Theorem~\ref{thm:gmap} below.

\begin{remA}[Numerical illustration]\label{rem:T2num}
With the constants of Remark~\ref{rem:T1num}, \eqref{eq:dhalf} gives $d_{1/2}=0.54$, and \eqref{eq:detectprob} gives leaf abstention probabilities $0.011$, $0.43$, $0.98$ and $1.00$ at true shifts $d=0$, $0.5$, $1$ and $2$. The limit \eqref{eq:pinf} at $d=0$ is $0.993$.
\end{remA}

\subsection*{C.2 The borrowing map $P(|g(x)|<c\given\text{data})$}

The map we display is the posterior probability that the local discrepancy is smaller than a pre-specified outcome-scale threshold $c>0$, $\mathcal B_c(x)=P(|g(x)|<c\given\text{data})$, together with the all-spike map $B(x)=P(\text{all }H_g\text{ leaves covering }x\text{ are spike}\given\text{data})$ as a diagnostic. For a single tree and a fixed partition, $g(x)=\theta_\ell$ for $x$ in leaf $\ell$, and the two are related by the leaf mixture of Theorem~\ref{thm:T1}(a).

\begin{thmA}[The discrepancy map for a single tree]\label{thm:gmap}
Under Assumption~\ref{ass:A1}, for a single $g$ tree with a fixed partition, $f$ and $(\sigma_1^2,\tau_0^2,\tau_1^2,w)$ fixed, and $x$ in a leaf with statistics $(\bar y,n)$:
\begin{enumerate}[label=(\alph*),leftmargin=2.2em]
\item (\emph{Exact form.}) $\mathcal B_c(x)=p\,\Psi_1(c)+(1-p)\,\Psi_0(c)$ with
\[
\Psi_j(c)=\Phi\Bigl(\frac{c-\rho_j\bar y}{\sqrt{\rho_jv}}\Bigr)-\Phi\Bigl(\frac{-c-\rho_j\bar y}{\sqrt{\rho_jv}}\Bigr),\qquad j\in\{0,1\}.
\]
\item (\emph{Relation to the spike probability.}) $\mathcal B_c(x)\le\Pspike(\bar y)+p\,\Psi_1(c)$, and on $\{\rho_0|\bar y|<c\}$,
\[
\mathcal B_c(x)\ \ge\ \Pspike(\bar y)\Bigl\{1-2\Phibar\Bigl(\frac{c-\rho_0|\bar y|}{\tau_0}\Bigr)\Bigr\}.
\]
\item (\emph{Consistency.}) If the RCT residuals in the leaf are $d+\epsilon_i$ as in Section~C.1 and $|d|\ne c$, then as $n\to\infty$, almost surely $\mathcal B_c(x)\to\mathbf 1\{|d|<c\}$, whatever the limit of $\Pspike$.
\end{enumerate}
\end{thmA}

\begin{proof}
(a) By Theorem~\ref{thm:T1}(a), the posterior of $g(x)=\theta$ is the two-component mixture with components $\Normal(\rho_j\bar y,\rho_jv)$, and $\Psi_j(c)$ is the probability that component $j$ lies in $(-c,c)$.

(b) The upper bound uses $\Psi_0\le1$. For the lower bound, $\mathcal B_c\ge(1-p)\Psi_0(c)$ and, for a $\Normal(m,V)$ variable with $|m|<c$, $P(|N|\ge c)\le P(|N-m|\ge c-|m|)=2\Phibar((c-|m|)/\sqrt V)$; here $|m|=\rho_0|\bar y|$ and $V=\rho_0v=\tau_0^2v/(\tau_0^2+v)\le\tau_0^2$, and $\Phibar$ is decreasing.

(c) Almost surely $\bar y\to d$, and $\rho_j\to1$, $\rho_jv\to0$ for both $j$, so $\rho_j\bar y\to d$ and each component concentrates at $d$: $\Psi_j(c)\to\mathbf 1\{|d|<c\}$ when $|d|\ne c$. A convex combination of two sequences with the same limit has that limit, regardless of the weights.
\end{proof}

\begin{remA}\label{rem:whygmap}
Parts (b) and (c) are the theoretical counterpart of the empirical finding that $\mathcal B_c$ separates shifted from compatible regions while $B(x)$ does not. The spike indicator answers whether the leaf discrepancy is of the order of $\tau_0$, a question whose answer stays uncertain in the limit (Theorem~\ref{thm:T2}(c)); $\mathcal B_c$ answers whether it is smaller than $c$ on the outcome scale, a question that the RCT controls settle as $n$ grows. The threshold $c$ must be pre-specified on the outcome scale; the study uses one third of $\sigma_{\rm RCT}$.
\end{remA}

\subsection*{C.3 The ensemble with fixed partitions, and what is not proved}

We return to the fixed-partition linear form \eqref{eq:linform}. For a profile $x$, let $e_x\in\{0,1\}^{L_g}$ select the $H_g$ leaves covering $x$, so $g(x)=e_x\tr\theta$. The atoms of the common refinement of the $H_g$ partitions are the cells $A_1,\dots,A_m$; $e_x$ is constant on each cell, and if $x$ and RCT control $i$ lie in the same cell then $e_x$ is row $i$ of $X$. Let $n_{1k}$ be the number of RCT controls in cell $A_k$.

\begin{propA}[The discrepancy map for $H_g$ trees]\label{prop:ensemble}
Under Assumption~\ref{ass:A1} and the conditioning of Section~B.2:
\begin{enumerate}[label=(\alph*),leftmargin=2.2em]
\item (\emph{Exact form.}) $\mathcal B_c(x)=\sum_zP(z\given r)\,\Psi_z(c)$ with $\Psi_z(c)=P\{|\Normal(e_x\tr\hat\theta(z),\,e_x\tr\Sigma_ze_x)|<c\}$, and $B(x)=\sum_{z\in\mathcal Z_x}P(z\given r)$ where $\mathcal Z_x=\{z:z_j=1\ \text{for all }j\text{ with }(e_x)_j=1\}$.
\item (\emph{Union bound.}) $1-B(x)\le\sum_{h=1}^{H_g}P(z_{h,\ell_h(x)}=0\given r)$ and $B(x)\le\min_hP(z_{h,\ell_h(x)}=1\given r)$, where $\ell_h(x)$ is the leaf of tree $h$ containing $x$. For $z\in\mathcal Z_x$, $e_x\tr\Sigma_ze_x\le e_x\tr D_ze_x=H_g\tau_0^2$, so on $\{|e_x\tr\hat\theta(z)|<c\}$, $\Psi_z(c)\ge1-2\Phibar\{(c-|e_x\tr\hat\theta(z)|)/(\sqrt{H_g}\tau_0)\}$.
\item (\emph{Consistency at fixed partitions.}) Suppose the true residual means are representable on the partitions, $\E[r\given x]=X\theta^\ast$ for some $\theta^\ast\in\mathbb R^{L_g}$, and let $x\in A_k$ with $m_k^\ast=e_x\tr\theta^\ast$ and $|m_k^\ast|\ne c$. If $n_{1k}\to\infty$ (the other cell counts arbitrary), then $\mathcal B_c(x)\to\mathbf 1\{|m^\ast_k|<c\}$ in probability.
\end{enumerate}
\end{propA}

\begin{proof}
(a) is \eqref{eqA:zpost} and the definition of the two maps. (b) The first two inequalities are the union bound and monotonicity for the events $\{z_{h,\ell_h(x)}=0\}$, $h=1,\dots,H_g$. For Gaussians the posterior covariance is dominated by the prior covariance, $\Sigma_z=(X\tr X/\sigma_1^2+D_z^{-1})^{-1}\preceq D_z$, and for $z\in\mathcal Z_x$, $e_x\tr D_ze_x=H_g\tau_0^2$; the tail bound is as in Theorem~\ref{thm:gmap}(b).

(c) Fix $z$; write $e=e_x$, $G=X\tr X/\sigma_1^2$, $m=e\tr\theta$. \emph{Variance.} The data enter the Gaussian model \eqref{eq:linform} through the cell means $\bar r_{k'}\given\theta\indep\Normal(e^{(k')\tr}\theta,\sigma_1^2/n_{1k'})$, where $e^{(k')}$ is the common row of $X$ in cell $k'$. In the joint Gaussian law of $(m,\bar r_1,\dots,\bar r_m)$ given $z$, conditional variances are constants and the law of total variance gives $\Var(m\given\bar r_k)=\E[\Var(m\given\bar r)\given\bar r_k]+\Var(\E[m\given\bar r]\given\bar r_k)\ge\Var(m\given\bar r)$; and $\Var(m\given\bar r_k)=\omega^2(\sigma_1^2/n_{1k})/(\omega^2+\sigma_1^2/n_{1k})\le\sigma_1^2/n_{1k}$ with $\omega^2=e\tr D_ze$. Hence $e\tr\Sigma_ze\le\sigma_1^2/n_{1k}\to0$. \emph{Mean.} With $r=X\theta^\ast+\epsilon$, \eqref{eqA:zpost} gives $\E[\theta\given z,r]=\Sigma_zG\theta^\ast+\Sigma_zX\tr\epsilon/\sigma_1^2$, and $\Sigma_zG=I-\Sigma_zD_z^{-1}$, so
\[
e\tr\E[\theta\given z,r]-m_k^\ast=-e\tr\Sigma_zD_z^{-1}\theta^\ast+e\tr\Sigma_zX\tr\epsilon/\sigma_1^2 .
\]
By the Cauchy--Schwarz inequality in the $\Sigma_z$ inner product,
\[
|e\tr\Sigma_zD_z^{-1}\theta^\ast|\le\sqrt{e\tr\Sigma_ze}\,\sqrt{(D_z^{-1}\theta^\ast)\tr\Sigma_z(D_z^{-1}\theta^\ast)}\le\sqrt{\sigma_1^2/n_{1k}}\;\tau_1\|\theta^\ast\|/\tau_0^2 ,
\]
using $\Sigma_z\preceq D_z\preceq\tau_1^2I$ and $\|D_z^{-1}\|\le\tau_0^{-2}$. The noise term has mean zero and variance $e\tr\Sigma_zG\Sigma_ze\le e\tr\Sigma_ze\le\sigma_1^2/n_{1k}$, because $G\preceq G+D_z^{-1}=\Sigma_z^{-1}$. Both terms vanish in probability. \emph{Conclusion.} For each of the finitely many $z$, $\Psi_z(c)=P\{|\Normal(\hat m_z,\hat V_z)|<c\}$ with $\hat m_z\to m_k^\ast$ in probability and $\hat V_z\to0$, so $\Psi_z(c)\to\mathbf 1\{|m_k^\ast|<c\}$ in probability when $|m_k^\ast|\ne c$, uniformly over $z$; the mixture $\sum_zP(z\given r)\Psi_z(c)$ then has the same limit whatever the weights.
\end{proof}

\paragraph{Scope under random partitions.}
Everything in this appendix about the map is conditional on the partitions of the $g$ trees. Three things are not proved. First, the representability assumption in Proposition~\ref{prop:ensemble}(c) is a real restriction: with two trees splitting on different variables the incidence matrix has rank three on four cells, so a discrepancy with an interaction on those cells is not representable, and the fitted cell means then converge to a weighted projection of the truth onto the additive span of the partitions rather than to the truth. Second, we have no statement about the posterior over the partitions themselves, hence none about $\mathcal B_c(x)$ or $B(x)$ under the full posterior; posterior concentration for a single standard BART ensemble is known \citep{rockova2020bart}, but the two-ensemble decomposition with a spike-and-slab discrepancy prior is outside the available theory, and the empirical separation of regions by $\mathcal B_c$ (Section~\ref{sec:sim}) is the evidence we have. Third, the leaf-level results treat $f$ as fixed; with $f$ estimated, the residual in a $g$ leaf carries the error of $f$ at the RCT profiles, which is $O_p$ of the RWD posterior standard deviation of $f$ there, and enters $d$ as an additive perturbation. On the estimand, this perturbation is what $V^f_\mu$ of Web Appendix~D measures.

\section*{Web Appendix D: Variance-based ESS and calibration (T3)}

We use a normal location experiment with known variance $\sigma_1^2$ as the reference information unit. For a smooth scalar prior density $p(\mu)$, the ELIR definition of \citet{neuenschwander2020} is
\[
\mathrm{ESS}_{\rm ELIR}=\sigma_1^2\E_p[-\partial_\mu^2\log p(\mu)].
\]
For a normal prior of variance $V$, this equals $\sigma_1^2/V$. For a mixture, the curvature is that of the log \emph{marginal} density, not the average of the component curvatures. Accordingly, our calibration is an average of working precision ratios, with $\sigma_1^2$ providing normal-reference units. It is not an exact marginal ELIR calculation and does not inherit ELIR's predictive-consistency property. In the survival analysis these units refer to uncensored log times.

\subsection*{D.1 Leaf aggregation and the calibration reference}

At fixed discrepancy partitions $\mathcal T$, define $a_{h\ell}=N_{h\ell}/N$ from the profiles used to standardize the estimand. With $\mu_g=\sum_{h,\ell}a_{h\ell}\theta_{h\ell}$, conditional independence of the leaf priors gives
\begin{equation}
V_g(\mathcal T,z,t)=\Var(\mu_g\given\mathcal T,z,\tau_0^2=t)
=\sum_{h,\ell}a_{h\ell}^2\{z_{h\ell}t+(1-z_{h\ell})\tau_1^2\}.
\label{eqA:leafvariance}
\end{equation}
Leaves with $a_{h\ell}=0$ contribute nothing. Let $L_+$ count the remaining leaves. For fixed $w$ and a specified scale law $P_s$, averaging the corresponding working precision over independent leaf states gives
\[
\mathrm{ESS}_{\mathcal T}(w,s)
=\sigma_1^2\sum_{z\in\{0,1\}^{L_+}}w^{|z|}(1-w)^{L_+-|z|}
 \int\frac{P_s(\dd t)}{V_\mu^f+V_g(\mathcal T,z,t)}.
\]
An additional expectation over $\mathcal T$ averages over the partition prior. When one leaf per tree has positive weight, its weight is one, $L_+=H_g$, and the sum reduces to a binomial average with variance $kt+(H_g-k)\tau_1^2$. This includes a pointwise profile and stump-only ensembles. With several occupied leaves per tree, their squared proportions and independent indicators must be retained. Merely counting the total number of spikes is insufficient when their weights differ.

For the informative calibration reference we set every discrepancy leaf to the spike. Then $V_g=tC(\mathcal T)$, where $C(\mathcal T)=\sum_{h,\ell}a_{h\ell}^2$ and $\E C(\mathcal T)=c_gH_g$. The implemented convention replaces $C(\mathcal T)$ by this mean and uses an untruncated scale law:
\begin{equation}
\ess_0(s_0^2)=\sigma_1^2\,\E_{\tau_0^2\sim\text{scaled-Inv-}\chi^2(\nu_0,s_0^2)}\Bigl[\frac1{V^f_\mu+c_gH_g\tau_0^2}\Bigr].
\label{eqA:ess0}
\end{equation}
Here $V_\mu^f>0$ is the variance of the standardized RWD-only BART surface. Its use represents that marginal distribution by its variance. By convexity, substituting $\E C$ before inversion gives no larger a value than averaging $1/(V_\mu^f+tC)$ over partitions at fixed $t$. Averaging these conditional precisions over $t$ also differs from taking the inverse of the variance after mixing over $t$. These are distinct operations.

For block variances $V^{(b)}$, the grid rule is
\begin{equation}
\hat s^2_{0\star}=\inf\{s_0^2\in\mathcal S:\widehat G(s_0^2)\ge q\},\qquad
\widehat G(s_0^2)=\frac1B\sum_{b=1}^B\mathbf 1\{\ess_0(V^{(b)};s_0^2)\le N_{\rm target}\}.
\label{eqA:prrule}
\end{equation}
Write $G(s)=P\{\ess_0(V;s)\le N_{\rm target}\}$ for a specified positive block-variance law. The theorem concerns exact evaluation of \eqref{eqA:ess0}, a fixed coefficient $c_g$, and the rule \eqref{eqA:prrule} before implementation overrides.

\begin{thmA}[Monotone, well-posed calibration]\label{thm:T3}
Fix $\sigma_1^2>0$, $V>0$, $c=c_gH_g>0$, $\nu_0>0$, $N_{\rm target}>0$ and $q\in(0,1)$, and write $s=s_0^2$.
\begin{enumerate}[label=(\alph*),leftmargin=2.2em]
\item (\emph{Monotonicity and limits.}) Under the untruncated law, $\ess_0(V;s)$ is continuous, differentiable, convex and strictly decreasing, with limits $\sigma_1^2/V$ and $0$ as $s\downarrow0$ and $s\uparrow\infty$. Under truncation to $(0,\tau_1^2)$ it remains continuous and strictly decreasing, with limits $\sigma_1^2/V$ and $\sigma_1^2/(V+c\tau_1^2)>0$. Convexity is asserted only for the untruncated law.
\item (\emph{Population selection.}) Under the untruncated law, extend $\ess_0(V;0)=\sigma_1^2/V$. Then $\{s\ge0:\ess_0(V;s)\le N_{\rm target}\}=[s^\dagger(V),\infty)$, where $s^\dagger$ is the unique positive root when $N_{\rm target}<\sigma_1^2/V$ and zero otherwise. The function $s^\dagger$ is continuous and strictly decreasing in $V$ where positive. Thus $G(s)=P\{s^\dagger(V)\le s\}$ is a distribution function, and its $q$-quantile $s_\star$ is the unique attained infimum of $\{G\ge q\}$. Moreover,
\begin{equation}
s_\star>0\quad\Longleftrightarrow\quad P\{N_{\rm target}<\sigma_1^2/V\}>1-q.
\label{eq:feasible}
\end{equation}
Under truncation, a positive finite root requires the target to lie strictly between the lower limit in (a) and the ceiling. At or below the lower limit, no finite scale satisfies the target.
\item (\emph{Monte Carlo limits.}) Let the full Stage-1 state process be a stationary, aperiodic Harris-ergodic Markov chain. Suppose the standardized surface is a square-integrable measurable function of that state, with variance $V_*>0$. Let $U_b$ be the sample variance in its $b$th block of length $m\ge2$, and set $V^{(b)}=\gamma_BU_b$ with $\gamma_B=\widehat V_{\rm whole}/(B^{-1}\sum_bU_b)$.
For fixed $B$ and $m\to\infty$, every $V^{(b)}\to V_*$ almost surely. On a finite grid avoiding equality $\ess_0(V_*;s)=N_{\rm target}$ and containing an admissible point, selection by \eqref{eqA:prrule} eventually equals the smallest admissible grid value.
For fixed $m$ and $B\to\infty$, suppose $U_b>0$ almost surely and $u_m=\E U_b>0$. Put $\gamma=V_*/u_m$ and
\[
G_m(s)=P\{\ess_0(\gamma U_b;s)\le N_{\rm target}\}.
\]
At any $s$ for which $P\{\ess_0(\gamma U_b;s)=N_{\rm target}\}=0$, $\widehat G(s)\to G_m(s)$ almost surely. On a finite grid satisfying this condition, with $G_m(s)\ne q$ at every point and at least one $G_m(s)>q$, the empirical and population grid selections eventually agree.
\item (\emph{Single-tree reference.}) Condition on one fixed $f$ tree and $\sigma_2^2$. Let $a_\ell$ and $\pi_\ell=n_{2\ell}/n_2$ be target and RWD leaf proportions, with $\pi_\ell>0$ whenever $a_\ell>0$. Terms with $a_\ell=\pi_\ell=0$ contribute zero to the sums below. Then
\[
V_\mu^f=\sum_\ell a_\ell^2(\lambda_f^{-2}+n_{2\ell}/\sigma_2^2)^{-1}
\le\frac{\sigma_2^2}{n_2}\sum_\ell\frac{a_\ell^2}{\pi_\ell},\qquad
\frac{\sigma_1^2}{V_\mu^f}\ge\frac{n_2\sigma_1^2/\sigma_2^2}{\sum_\ell a_\ell^2/\pi_\ell}.
\]
Equality is attained in the flat-prior limit. In that limit the denominator is at least one, with equality exactly when $a=\pi$. At finite $\lambda_f^2$ the prior contributes information, and covariate shift need not lower the ceiling relative to its value at $a=\pi$.
\end{enumerate}
\end{thmA}

\begin{proof}
(a) Write $\tau_0^2=\nu_0s/Q$, $Q\sim\chi^2_{\nu_0}$. The integrand $(V+c\nu_0s/Q)^{-1}$ is continuous, strictly decreasing and convex in $s$, bounded by $1/V$. Dominated convergence gives continuity and the limits; integration preserves monotonicity and convexity. Its derivative is bounded in absolute value by $1/(4Vs)$, which permits differentiation on compact subsets of $(0,\infty)$. Under truncation, the density ratio at $s'>s$ is proportional to $\exp\{-\nu_0(s'-s)/(2t)\}$, strictly increasing in $t$. The truncated family is therefore stochastically increasing, giving strict decrease of the expected reciprocal. Continuity follows from integration on a fixed interval. As $s\downarrow0$, the untruncated law concentrates at zero and its truncation probability tends to one. As $s\uparrow\infty$, compare any interval below $\tau_1^2-\epsilon$ with one inside $(\tau_1^2-\epsilon/2,\tau_1^2)$: the exponential density factor makes the former probability tend to zero. The truncated law concentrates at $\tau_1^2$, giving the stated lower limit.

(b) Continuity, strict decrease and the limits in (a) give the root and sublevel set. Joint continuity and strict decrease in $V$ give continuity and strict decrease of a positive root. Since $s^\dagger(V)$ is finite almost surely, $G$ is a distribution function with an attained finite $q$-quantile. Right-continuity gives $s_\star>0$ exactly when $G(0)<q$, equivalent to \eqref{eq:feasible}. The truncated range in (a) gives its feasibility restriction.

(c) At fixed $B$, ergodic laws for the draw and its square imply convergence of every growing block variance to $V_*$. A block sum is the difference of two partial sums, so the same conclusion holds for later blocks whose starting index also increases. The whole-chain and mean block variances have that limit, giving $\gamma_B\to1$. Continuity and strict grid separation yield the selection result.
At fixed $m$, the nonoverlapping block sequence is stationary and ergodic. Thus $B^{-1}\sum_bU_b\to u_m$ and $\gamma_B\to\gamma$ almost surely. For each small $\epsilon>0$, eventually $(\gamma-\epsilon)U_b\le\gamma_BU_b\le(\gamma+\epsilon)U_b$ for every block. Monotonicity in $V$ bounds the empirical indicators between those at the two fixed scale factors. Apply the ergodic theorem to each bound and let $\epsilon\downarrow0$; absence of mass at the threshold makes both limits $G_m(s)$. A finite set of grid points with $G_m(s)\ne q$ eventually has the correct sign relative to $q$, proving selection. No central limit theorem or convergence rate is asserted.

(d) The fixed-tree posterior leaf variances are $(\lambda_f^{-2}+n_{2\ell}/\sigma_2^2)^{-1}$ and the leaves are independent, giving the expression for $V_\mu^f$. Drop $\lambda_f^{-2}$ for the upper bound. Cauchy--Schwarz gives $1=(\sum_\ell a_\ell)^2\le\sum_\ell a_\ell^2/\pi_\ell$, with equality exactly at $a=\pi$. If a target leaf has no RWD observations, its finite-prior variance remains $\lambda_f^2$; the flat-prior reference for that target is unavailable.
\end{proof}

\begin{remA}[Units and regularization]\label{rem:units}
The fitted ceiling $\sigma_1^2/V_\mu^f$ is a property of the working precision summary. Both external observations and the BART prior determine $V_\mu^f$. The single-tree reference in part (d) does not justify subtracting a universal prior contribution from a fitted ensemble. Under the range rule, $H_f\lambda_f^2=(\max y-\min y)^2/(4k^2)$, so increasing the number of trees alone does not remove prior information.
\end{remA}

\begin{remark}[Calibration implementation]
Stage 1 fits BART to the RWD and evaluates the standardized surface at the target profiles. In two-arm analyses, $\sigma_1^2$ is the posterior mean residual variance from a separate BART fitted to the RCT controls; a design value or the RWD estimate is used when controls are unavailable. The code fixes $B=20$, estimates $c_g$ with $2{,}000$ prior trees, and evaluates the untruncated scale expectation with $4{,}000$ common chi-squared draws across the grid. These finite approximations are separate from the exact-integral theorem. For a fixed positive variance and finite grid, their errors vanish as the inner sample sizes grow; conclusions about selection require strict separation at the limiting grid values.

A requested target at or above the whole-chain ceiling triggers an override: the smallest grid scale is returned, without requiring the block rule to pass. Otherwise the first grid value satisfying the rule is selected. If none passes, the largest value is returned with an upper-boundary flag; it does not thereby meet the target. A small scale approximates complete pooling of $f_1$ with $f$ only when $w=1$. With $w<1$, slab discrepancies remain possible. Equation~\eqref{eqA:ess0} uses an untruncated calibration law, while posterior fitting uses \eqref{eq:tau0prior}; replacing either law or the boundary rule would change the fitted analyses.

The realized diagnostic averages $\sigma_1^2/\{V_\mu^f+V_g(\mathcal T,z,\tau_0^2)\}$ over posterior states, using \eqref{eqA:leafvariance}. This evaluates prior variance at learned states; it is not the posterior variance of the discrepancy or a subtraction of posterior precisions. The application map uses $\hat\sigma_1^2/\{V_f(x)+H_g\bar\tau_0^2\}$, with the fitted mean scale $\bar\tau_0^2$. In single-arm mode this estimates the mean under the calibrated truncated prior. The supplementary pointwise prior summaries instead average $\hat\sigma_1^2/\{V_f(x)+H_g\tau_0^2\}$ over the untruncated calibration prior. The inversion and expectation occur in different orders. All regional summaries use their own standardization profiles and are not additive.
\end{remark}

\section*{Web Appendix E: Identifiability of $f$ and $g$}

\begin{propA}[Population identification]\label{prop:ident}
Let $\mathcal X_1$ and $\mathcal X_2$ be the supports of the covariate distributions of the RCT controls and of the RWD, and suppose the regression functions $x\mapsto\E[y\given x,D=0]$ on $\mathcal X_2$ and $x\mapsto\E[y\given x,D=1]$ on $\mathcal X_1$ are known (the infinite-data limit). Under model \eqref{eqA:model}:
\begin{enumerate}[label=(\alph*),leftmargin=2.2em]
\item $f$ is identified on $\mathcal X_2$ and $f+g$ is identified on $\mathcal X_1$.
\item $g$ is identified on $\mathcal X_1\cap\mathcal X_2$, as $g=(f+g)-f$.
\item On $\mathcal X_1\setminus\mathcal X_2$ only the sum $f+g$ is identified; the split between $f$ and $g$ is determined by the prior, with $f$ the BART extrapolation of the RWD surface and $g$ the remainder.
\item On $\mathcal X_2\setminus\mathcal X_1$, $f$ is identified but $g$ is not identified from the source-specific outcome distributions. Its posterior predictions depend on extrapolation through tree partitions and on the prior. This population statement does not imply that its fitted posterior law at those profiles equals the prior.
\end{enumerate}
These statements concern unrestricted population regression functions. Fixed tree partitions impose additional restrictions that can transmit information across these regions. In the single-arm case $\mathcal X_1=\emptyset$, the likelihood does not involve $g$ anywhere. Its entire posterior law equals its calibrated prior, and, conditional on the scales and $w$, its pointwise variance is $H_g\{w\tau_0^2+(1-w)\tau_1^2\}$.
\end{propA}

\begin{proof}
(a) is the model: $\E[y\given x,D=0]=f(x)$ on $\mathcal X_2$ and $\E[y\given x,D=1]=f(x)+g(x)$ on $\mathcal X_1$. (b), (c), (d) are immediate from (a) and from the fact that no other moment of the data involves $f$ or $g$. For the single-arm variance, condition on the scales and $w$. Then $g(x)=\sum_h\theta_{h,\ell_h(x)}$ is a sum of $H_g$ independent leaf contributions, each with variance $w\tau_0^2+(1-w)\tau_1^2$.
\end{proof}

\begin{remA}[Finite samples]\label{rem:identfinite}
In finite samples the statement is leaf-wise. A $g$ leaf with no RCT controls has $M_g=1$ (Lemma~\ref{lem:marg}(b)) and its $(z,\theta)$ is drawn from the prior at the current $(\tau_0^2,w)$; a $g$ leaf that straddles the RCT support extends the RCT-informed discrepancy as a constant into the region without RCT controls, which is a piecewise-constant extrapolation and not a prior draw. On $\mathcal X_1\cap\mathcal X_2$ the separation of $f$ and $g$ rests on the two priors: $f$ is pinned by the RWD, whose observations do not involve $g$, and the spike shrinks $g$ toward zero, so a shift common to both sources is attributed to $f$ and a shift of the RCT controls alone to $g$. The mixture \eqref{eqA:leafprior} is identifiable as a two-component scale mixture as long as $\tau_0^2\ne\tau_1^2$, and Assumption~\ref{ass:A1} fixes which component is the spike, so the label of $z$ is well defined.
\end{remA}

\section*{Web Appendix F: Scope of the posterior-invariance result}

The leaf marginals in Lemma~\ref{lem:marg} and the marginal-then-conditional update in Proposition~\ref{prop:exact} establish invariance for the specified posterior kernel. This construction uses the global discrepancy scale as a conditioned-on parameter during a tree move and updates it separately from its truncated conditional. Thus no moment-matched likelihood or approximation-error bound enters the tree acceptance ratio. Invariance does not imply independent posterior draws, rapid mixing, or adequate exploration by a finite run. Those properties require separate analysis and diagnostics. Likewise, the conditional influence and map results in Appendices B and C and the working calibration result in Appendix D have their own assumptions; the posterior-invariance proposition does not establish them for the full fitted ensemble.
